\documentclass[11pt]{article}

\usepackage[margin=1in]{geometry}
\usepackage[T1]{fontenc}
\usepackage{lmodern}
\usepackage{amsmath,amssymb,mathtools,bm}
\usepackage{booktabs}
\usepackage{enumitem}
\usepackage{xcolor}
\usepackage{microtype}
\usepackage[hidelinks]{hyperref}

\definecolor{navy}{RGB}{25,55,92}
\definecolor{muted}{RGB}{90,96,105}
\newcommand{\R}{\mathbb{R}}
\newcommand{\E}{\mathbb{E}}
\newcommand{\Prob}{\mathbb{P}}
\newcommand{\Ind}{\mathbf{1}}
\newcommand{\simplex}{\Delta}
\newcommand{\argmax}{\operatorname*{arg\,max}}
\newcommand{\argmin}{\operatorname*{arg\,min}}
\newcommand{\Dip}{\operatorname{Dip}}
\newcommand{\diag}{\operatorname{diag}}
\newcommand{\nullgen}{\widehat{\mathcal P}_0}
\newcommand{\why}[1]{\par\noindent\textcolor{muted}{\small\itshape Why this step: #1}}

\hypersetup{
  pdftitle={Detect, Recover, Interpret (v2): Calibrated Black-Box Auditing of Hidden Routing Regimes},
  pdfsubject={Revised working algorithm note},
  pdfkeywords={black-box auditing, regime detection, null generator, selection-aware bootstrap, label synchronization, certified propagation}
}

\title{\textcolor{navy}{Detect, Recover, Interpret:\\
Black-Box Auditing of Hidden Routing Regimes\\
\large Revised algorithm note (v2): calibrated detection and certified recovery}}
\author{Working algorithm note}
\date{\today}

\begin{document}
\maketitle

\begin{abstract}
This revision restates the detection target and hardens the machinery of the three-stage
auditing framework. The principal changes relative to v1 are:
(1) the Stage~A null is reframed --- the tested hypothesis is compatibility with a
\emph{calibrated smooth single-branch reference distribution}, not the literal statement
$K=1$; rejection is evidence of boundary-like multiscale response geometry, consistent with
latent routing subject to corroboration;
(2) the ratio statistic $R_i$ is replaced by a stabilized log-range statistic and, preferably,
a studentized deviation from a fitted smooth-null scale curve, with a minimum-signal
abstention rule;
(3) the null generator $\nullgen$ is made concrete as a hierarchical Gaussian-process model of
single-branch local response surfaces, fitted robustly and cross-fitted for testing;
(4) Benjamini--Hochberg screening is demoted to a screening device; the primary confirmatory
claim is an omnibus selection-aware parametric bootstrap whose $p$-value has standard
finite-simulation form and whose calibration is conditional on generator adequacy;
(5) shape corroboration specifies its scalar projections explicitly;
(6) Stage~B uses whitened features and soft mixture responsibilities, commits to a concrete
partial-matching label-synchronization algorithm, fits gates by regularized responsibility-weighted
logistic regression with stage-repeating bootstrap uncertainty, and replaces propagation
through unflagged regions with bridge-probe \emph{edge certification}, on the principle that
non-rejection is not evidence of interiority;
(7) the audit design adds budget splitting, negative controls, and probe-support diagnostics;
(8) Stage~C interpretation language is weakened to what the evidence supports. A closing
section states what Stage~B still requires to reach Stage~A's level of maturity.
\end{abstract}

\section{Problem, target, and the precise detection claim}

Let a black-box system be represented as
\[
  f(x)=f_{g(x)}(x)+\eta,
\]
where \(g:\mathcal X\to[K]\) is an unobserved routing rule, \(f_k\) is the response mechanism
in regime \(k\), and \(\eta\) is observation or sampling noise. The auditor observes only
queries and responses. For anchors \(x_1,\ldots,x_n\), a probe family
\(Q_\sigma(\cdot\mid x_i)\) generates perturbations
\[
  \delta_{ij}^{(t)}\sim Q_{\sigma_t}(\cdot\mid x_i),
  \qquad
  \psi_{ij}^{(t)}=\psi\!\left(f(x_i+\delta_{ij}^{(t)})\right)\in\R^V,
\]
with response feature map \(\psi:\mathcal Y\to\R^V\). The local routing composition is
\(\pi_i^\sigma(k)=\Prob_{\delta\sim Q_\sigma}[g(x_i+\delta)=k]\in\simplex^{K-1}\).
The three goals remain separated: \textbf{detection} (Stage~A), \textbf{recovery} (Stage~B),
and \textbf{interpretation} (Stage~C).

\subsection{What Stage A does and does not test}
\label{sec:claim}

A single smooth mechanism can produce curvature-driven scale dependence, spatially varying
derivatives, heteroskedastic output noise, anisotropic noise, sharp-but-continuous
transitions, and unstable normalization at locally flat anchors. Conversely, genuine routing
can remain invisible when branches overlap under \(\psi\), when probes move parallel to the
gate, or when the local mixture is effectively unimodal. Stage~A therefore does not, and
cannot, test the literal hypothesis \(K=1\). The tested null is
\[
  H_{0,i}:\ \text{anchor \(i\) follows the calibrated smooth single-branch null } \nullgen ,
\]
and the operational claim is:
\begin{quote}
\emph{Given a specified response representation \(\psi\), probe design
\(\{Q_{\sigma_t}\}\), and null class, identify anchors whose multiscale response geometry is
incompatible with the fitted reference distribution of single-branch smooth behavior.}
\end{quote}
Rejection is \emph{consistent with} boundary crossing or latent routing and must be
corroborated by shape evidence (Stage~A.5), recovery structure (Stage~B), and mechanism
analysis (Stage~C) before any routing language enters the audit finding. Both positive and
negative findings are relative to \(\psi\), the probe design, and the null class: a
sharp-but-continuous transition narrower than the probed scales is observationally equivalent
to a routing boundary, and a mechanism switch invisible through \(\psi\) is undetectable in
principle.

\section{Audit design and budget discipline}
\label{sec:design}

\subsection{Budget split}
\label{sec:budget}

Partition the query budget into four disjoint samples:
\begin{itemize}[leftmargin=1.8em]
  \item \(\mathcal S_1\) --- detection and scale selection (Stage A.1--A.4);
  \item \(\mathcal S_2\) --- modality confirmation at flagged anchors (Stage A.5);
  \item \(\mathcal S_3\) --- mixture recovery, synchronization, and gate fitting (Stage B);
  \item \(\mathcal S_4\) --- final held-out validation of every reported quantity.
\end{itemize}
The selection-aware bootstrap can in principle account for sample reuse, but a genuine
held-out split is easier to defend and removes the dependence of confirmatory claims on
correctly modeling the reuse. When the budget forces reuse, the bootstrap must reproduce it
exactly (\S\ref{sec:bootstrap}).

\subsection{Probe families and support diagnostics}
\label{sec:support}

Maintain at least two probe families: \emph{on-manifold} (perturbations consistent with the
data distribution) and \emph{ambient} (e.g., isotropic Gaussian in input space), across a
ladder \(\sigma_1<\cdots<\sigma_T\); feature-targeted probes may be added when a candidate
gating feature exists. For every anchor and probe family, report a \emph{support realism
score} for the perturbed inputs --- e.g., the quantile of their \(k\)NN distance to a
reference sample of real inputs. Findings are reported stratified by support realism:
``ambient-only'' evidence obtained at unrealistic inputs is explicitly qualified, since it may
reflect behavior in operationally meaningless regions.

\subsection{Negative controls}
\label{sec:controls}

Every audit run includes controls probed with the identical pipeline:
smooth curved single-regime functions; heteroskedastic single-regime functions;
discontinuous functions without latent routing (e.g., a single mechanism containing a hard
kink); known benign mixtures of experts; and randomized or irrelevant probe directions.
Controls establish whether the statistic responds specifically to boundary-like switching
rather than to generic nonlinearity or noise structure. The discontinuous control also makes
the observational-equivalence point of \S\ref{sec:claim} concrete: an intrinsic step inside
one mechanism \emph{will} be flagged, because at the level of observable behavior it is a
regime boundary; attributing mechanism is Stage~C's job, not Stage~A's.

\section{Stage A: multi-scale detection}

\subsection{A.1: Noise model, whitening, and the dispersion curve}
\label{sec:whiten}

Repeated unperturbed queries at each anchor (and pooled across anchors) yield an estimated
observation-noise covariance \(\widehat\Sigma_{\mathrm{obs}}\in\R^{V\times V}\); when \(V\) is
large relative to the number of repeats, use a shrinkage or low-rank estimator
(e.g., Ledoit--Wolf). Whiten all response features:
\[
  \tilde\psi_{ij}^{(t)}=\widehat\Sigma_{\mathrm{obs}}^{-1/2}\,\psi_{ij}^{(t)} .
\]
\why{scalar noise-floor subtraction implicitly assumed isotropic noise; whitening makes the
floor isotropic by construction, at the cost of requiring a stable covariance estimate.}

With \(\tilde S_i(\sigma_t)\) the empirical scatter of
\(\{\tilde\psi_{ij}^{(t)}\}_{j=1}^m\), the repeated-query scatter is \(\approx I\), so the
excess top-eigenvalue spread per unit input scale is
\[
  r_i(\sigma_t)
  =
  \frac{\sqrt{\bigl(\lambda_{\max}(\tilde S_i(\sigma_t))-1\bigr)_+}}{\sigma_t}.
\]
For a locally linear single branch with whitened directional slope \(b_i\),
\(\lambda_{\max}(\tilde S_i(\sigma))\approx 1+\sigma^2\|b_i\|^2\) and
\(r_i(\sigma)\approx\|b_i\|\); curvature makes \(r_i\) rise with \(\sigma\) even under the
null, which is why the null model must represent curvature (\S\ref{sec:generator}) rather
than the statistic assuming linearity.

\paragraph{Minimum-signal abstention.}
If \(\lambda_{\max}(\tilde S_i(\sigma_t))-1\) is not significantly positive at any scale
(one-sided test against the repeated-query eigenvalue distribution), the anchor is reported as
\emph{insufficient signal}, which is distinct from \emph{null-consistent}: a locally flat
model provides no information about boundaries, and a noise-dominated denominator must not
enter any ratio or range statistic.
\why{the v1 ratio could explode when one scale's spread was noise-underestimated; the anchor
looked anomalous for reasons unrelated to any bump.}

\subsection{A.2: Stabilized scale statistics}
\label{sec:statistic}

The v1 ratio \(R_i=\max_t r_i/\min_t r_i\) is deprecated: its denominator is unstable near
flat regions, and its null distribution conflates bump detection with normalization noise.
Two replacements, in increasing order of preference:
\begin{enumerate}[label=\textbf{S\arabic*.},leftmargin=2.6em]
  \item \textbf{Log-range.}
  \[
    R_i^{\log}
    =\max_t \log\bigl(r_i(\sigma_t)+\epsilon\bigr)
    -\min_t \log\bigl(r_i(\sigma_t)+\epsilon\bigr),
  \]
  with a fixed small \(\epsilon>0\), computed only on anchors passing the minimum-signal
  condition. A robust default when no null generator is yet fitted.
  \item \textbf{Studentized deviation from the fitted null scale curve.} With the null
  generator of \S\ref{sec:generator} fitted, obtain by simulation (or GP analytics) the null
  mean and standard deviation of the log dispersion curve at anchor \(i\),
  \(\mu_{0,i}(\sigma_t)\) and \(s_{0,i}(\sigma_t)\), and set
  \[
    T_i^{\mathrm{scale}}
    =\max_t\,
    \frac{\log r_i(\sigma_t)-\mu_{0,i}(\sigma_t)}{s_{0,i}(\sigma_t)},
    \qquad
    \sigma_i^\star=\argmax_t\,
    \frac{\log r_i(\sigma_t)-\mu_{0,i}(\sigma_t)}{s_{0,i}(\sigma_t)} .
  \]
  This asks the calibrated question directly: \emph{is the observed scale dependence larger
  than the fitted population of smooth single-branch behaviors predicts at this anchor?}
\end{enumerate}
\why{a bump statistic should measure deviation from the null's own scale dependence,
not deviation from flatness; curvature is a null feature, not evidence.}

\subsection{A.3: The null generator \texorpdfstring{$\nullgen$}{P0}}
\label{sec:generator}

The null generator is the object on which all calibrated inference rests, so it is specified
concretely. \(\nullgen\) is a fitted joint generative law over a \emph{complete synthetic
audit dataset} --- all responses at all anchors, scales, and probes, plus repeats --- under
the single-branch hypothesis. It must represent: (a) smooth local response surfaces with
realistic slopes and curvatures; (b) cross-anchor variation in those slopes and curvatures;
(c) observation noise, including heteroskedasticity across anchors; (d) spatial or batch
dependence among anchors.

\paragraph{Hierarchical GP specification (default).}
At anchor \(i\), model each whitened response coordinate (or each of the leading response
principal components) of the local surface as a Gaussian process in the perturbation:
\[
  m_i(\delta)\sim\mathcal{GP}\!\left(
    \beta_i^\top\delta,\;
    s_i^2\,k_\nu\!\left(\|\delta-\delta'\|/\ell_i\right)
  \right),
  \qquad
  \tilde\psi_{ij}\mid m_i \sim
  \mathcal N\!\left(m_i(\delta_{ij}),\,\Sigma_i^{\mathrm{noise}}\right),
\]
with Mat\'ern kernel \(k_\nu\). The hyperparameters have direct null meaning:
\(\beta_i\) is the local linear slope, \(s_i\) the amplitude of smooth nonlinearity,
\(\ell_i\) the curvature lengthscale, and \(\nu\) the differentiability class ---
``how smooth is smooth'' is an explicit modeling choice rather than an implicit one.
Cross-anchor structure is hierarchical: \((\log s_i,\log\ell_i,\beta_i)\) are drawn from
population laws fitted across anchors, with spatially smooth mean fields when anchors are
spatial; \(\Sigma_i^{\mathrm{noise}}\) comes from local repeats or a fitted variance
function. Simulation is direct: draw hyperparameters, draw the surface values at the actual
probe locations (a finite-dimensional Gaussian with the kernel Gram matrix), add noise ---
one synthetic audit per replicate, at the original anchors, ladder, and probe counts, with
spatial dependence simulated jointly by blocks.

\paragraph{Fitting under contamination.}
The null must be fitted from data that may contain alternatives. Options, in increasing order
of strength:
trimmed fitting that excludes anchors in the top fraction of a preliminary detection statistic
(report sensitivity to the trimming fraction);
a separate calibration anchor set never used for testing;
and \textbf{cross-fitting} --- partition anchors into folds, test each fold against a
generator fitted on the remaining folds. Cross-fitting is the default for anchor-level
screening $p$-values.
\why{fitting \(\widehat F_0\) or \(\nullgen\) on the same statistics being tested makes the
null random, contaminated, and dependent on every observation; empirical-null BH over such
fits has no automatic FDR guarantee.}

\paragraph{Misspecification honesty.}
The generator's smoothness class \emph{defines} the null of \S\ref{sec:claim}. A single
mechanism with a transition sharper than \(\ell_i\) at the probed scales is outside the null
and will be flagged; that is the intended behavior, and the finding is stated as boundary-like
geometry, not as proof of a latent gate. Calibration of every downstream $p$-value is
conditional on generator adequacy; the evaluation plan (\S\ref{sec:eval}) therefore includes
misspecification experiments (heavy-tailed noise, unmodeled heteroskedasticity, kinks,
wrong \(\nu\)) measuring realized error rates when the generator is wrong. Alternative null
classes (local polynomial with random coefficients, splines) are acceptable substitutes; the
GP default is chosen because smoothness is the null hypothesis and the GP parameterizes
exactly smoothness with few, interpretable, estimable hyperparameters.

\subsection{A.4: Screening and multiplicity}
\label{sec:screening}

Compute anchor-level screening $p$-values \(p_i\) for \(T_i^{\mathrm{scale}}\) (or
\(R_i^{\log}\)) by simulation from the cross-fitted generator. Apply Benjamini--Hochberg at
rate \(q\) to obtain the flagged set \(F\). In this framework BH is a \emph{screening
device} that concentrates the confirmation budget; it is not reported as a final FDR
guarantee, because the fitted null is estimated and anchors may be dependent. Dependence is
addressed by blockwise simulation in the generator, block-level testing, or
Benjamini--Yekutieli as sensitivity analyses. Formal error control lives in the omnibus
bootstrap claim (\S\ref{sec:bootstrap}); anchor-level confirmatory claims additionally require
the per-anchor selection-aware $p$-values with a dependence-robust multiplicity correction.
\why{screening and confirmation are different jobs; conflating them forced v1's empirical-null
BH to carry a calibration burden it could not meet.}

\subsection{A.5: Shape corroboration with specified projections}
\label{sec:dip}

A dispersion bump is corroborated by evidence of multiple response populations. Hartigan's
dip test is univariate, so the scalar projection must be specified:
\begin{enumerate}[label=\textbf{P\arabic*.},leftmargin=2.6em]
  \item \textbf{Default:} dip on the first principal-component score of the whitened
  confirmation sample (\(\mathcal S_2\)) at \(\sigma_i^\star\).
  \item \textbf{Prespecified set:} dip on a small fixed set of directions --- top two PCs
  plus, when available, a candidate gating direction --- with Bonferroni correction over the
  set.
  \item \textbf{Projection pursuit (optional):} maximize a bimodality index over directions;
  if used, the entire pursuit is repeated inside every bootstrap replicate.
\end{enumerate}
\paragraph{Known power limitation.} Principal components maximize variance, not between-regime
separation; strong within-regime anisotropy can hide a mixture from the top PC. A
non-significant dip therefore does not certify unimodality. Multivariate modality or
clustering-based tests may replace the dip at the cost of harder calibration.
\why{v1 said ``project onto leading directions and dip,'' which is ambiguous for \(d_i>1\)
and silently favorable if the direction is tuned on the same data.}

\subsection{A.6: Selection-aware bootstrap confirmation}
\label{sec:bootstrap}

When the same sample informs anchor, scale, or projection choices, confirmation must
reproduce the complete search. Define
\[
  T_i^{\mathrm{obs}}
  =\Ind\{i\in F^{\mathrm{obs}}\}\;
  \Dip\!\left\{u_i^{\mathrm{obs}\top}\tilde\psi_{ij}^{(\star,\mathrm{obs})}\right\}_{j=1}^m,
\]
with \(u_i^{\mathrm{obs}}\) the specified projection of \S\ref{sec:dip}; setting the
statistic to zero for unselected anchors makes screening part of the statistic. For
\(b=1,\ldots,B\): generate a complete synthetic audit from \(\nullgen\); re-estimate the noise
model and whitening; recompute dispersion curves and scale statistics; refit the cross-fitted
nulls; rerun screening to get \(F^{(b)}\); re-derive projections; recompute
\(T_i^{(b)}\). Every data-adaptive choice --- trimming fractions, fold assignments,
\(K\)-selection rules, projection pursuit if used --- is repeated inside each replicate.
The anchor-level bootstrap value is
\[
  \widehat p_i^{\mathrm{sel}}
  =\frac{1+\sum_{b=1}^B\Ind\{T_i^{(b)}\ge T_i^{\mathrm{obs}}\}}{B+1},
\]
which has the correct finite-simulation form (and cannot return zero); its calibration is
\emph{conditional on the adequacy of the fitted null generator} --- this is an approximate
parametric bootstrap, not an exact test. The \textbf{primary existence claim} uses the
omnibus statistics
\[
  T_{\max}^{\mathrm{obs}}=\max_i T_i^{\mathrm{obs}},
  \qquad
  T_{\max}^{(b)}=\max_i T_i^{(b)},
\]
whose tail probability automatically absorbs the search over anchors and avoids asserting
anchor-level FDR before the null machinery is mature. Resampling only selected projections,
or holding \(\sigma_i^\star\), \(u_i\), or \(F\) fixed after learning them, does not correct
selection. A held-out \(\mathcal S_2\)/\(\mathcal S_4\) confirmation remains the cleanest
design when the budget allows.

\section{Stage B.1: query-level local decomposition}

Stage B runs only at confirmed anchors, on the recovery sample \(\mathcal S_3\), in whitened
coordinates. Conditional on its perturbation, query \(j\) has one latent label
\(a_{ij}=g(x_i+\delta_{ij})\), even though the profile is an admixture.

\subsection{Local mixtures and soft responsibilities}
\label{sec:soft}

At \(i\in F\), project the recovery sample to \(d_i\) leading whitened directions and fit
mixtures for \(K_i^{\mathrm{loc}}\in\{1,\ldots,K_{\max}\}\), selected by held-out likelihood,
bootstrap cluster stability, a complexity criterion, and a minimum-mass rule. Retain the
posterior responsibilities
\[
  \gamma_{ij}(\ell)
  =\widehat{\Prob}\bigl(a_{ij}=\ell\mid\tilde\psi_{ij}\bigr),
\]
component centroids \(\widehat\mu_{i\ell}\) and covariances \(\widehat C_{i\ell}\), and
estimate compositions \emph{softly}:
\[
  \widehat\pi_i(\ell)=\frac{1}{m}\sum_{j=1}^m\gamma_{ij}(\ell).
\]
\why{hard-assignment frequencies are biased toward \(1/K\) under component overlap;
responsibilities are the consistent estimator and propagate assignment uncertainty
downstream.}
The minimum-mass rule uses soft mass, \(m\,\widehat\pi_i(\ell)\ge n_{\min}\). The unperturbed
response \(\tilde\psi_{i0}\) is assigned by posterior to a local component
(\(\widehat z_i^{\mathrm{loc}}\)), identifying which branch owns the anchor; it is never
treated as a pure global archetype.

\section{Stage B.2: local-to-global alignment as partial-permutation synchronization}
\label{sec:sync}

Local labels are arbitrary, and --- because not every regime appears in every neighborhood ---
the alignment problem is synchronization over \emph{partial} matchings, which is harder than
synchronization over complete permutations. The committed algorithm:

\begin{enumerate}[label=\textbf{SYNC\arabic*.},leftmargin=3.4em]
  \item \textbf{Graph.} Build the anchor graph \(G\) from geographic, manifold, or
  task-metric neighborhoods.
  \item \textbf{Edge costs.} For edge \((i,j)\) and components \((\ell,r)\),
  \[
    c_{ij}(\ell,r)
    =(\widehat\mu_{i\ell}-\widehat\mu_{jr})^\top
     (\widehat C_{i\ell}+\widehat C_{jr}+\lambda I)^{-1}
     (\widehat\mu_{i\ell}-\widehat\mu_{jr}),
  \]
  optionally augmented with local-slope and held-out predictive terms.
  \item \textbf{Partial matching.} Solve a rectangular assignment with an unmatched
  (dummy) option of cost \(\chi\), so components may remain unmatched. Default
  \(\chi\): a high quantile (e.g.\ 0.95) of same-branch matching costs simulated from the
  fitted within-regime continuity model, with sensitivity reported.
  \item \textbf{Edge confidence.} Confidence of edge \((i,j)\) is the margin between its best
  and second-best total assignment cost.
  \item \textbf{Spanning forest and propagation.} Take the maximum-confidence spanning forest
  of \(G\) (Kruskal on edge confidence) and propagate labels along it to define provisional
  global classes.
  \item \textbf{Cycle audit and abstention.} Every off-forest edge is re-matched
  independently; the component-level cycle-consistency rate \(\rho_{cc}\) is the
  confidence-weighted fraction of off-forest edges agreeing with the forest labeling.
  Components with \(\rho_{cc}\) below threshold are abstained from, not forcibly
  synchronized.
  \item \textbf{Global count.} \(\widehat K\) is the number of global classes whose aggregate
  soft mass exceeds a floor \emph{and} whose class survives in a stated fraction of Stage-B
  bootstrap replicates (\S\ref{sec:gateuq}).
\end{enumerate}
\why{v1 said ``synchronize to cycle consistency,'' a problem statement rather than a
procedure; the forest-plus-audit construction is implementable, and its failure mode
(low \(\rho_{cc}\)) is an explicit abstention trigger.}

Raw global clustering remains rejected: within one regime the response varies across the
input space (\(w_k(x_i)\neq w_k(x_j)\)), so global clusters follow magnitude and covariates
rather than mechanism; the local-to-global route follows smooth branches through overlapping
neighborhoods. The three-regime Slack-style stack illustrates assembly: local views
\(\{1,2\}\), \(\{1,3\}\), \(\{2,3\}\) under different probe families chain into one
three-regime system through shared components, even if no unperturbed anchor sits in regime 3.

\section{Stage B.3: gates, uncertainty, and certified propagation}

\subsection{Regularized, responsibility-weighted gate fits}
\label{sec:gate}

At a flagged anchor with a locally binary boundary, fit
\[
  \Prob(a_{ij}=\ell\mid\delta_{ij})
  =\operatorname{logit}^{-1}(\beta_{i0}+\beta_i^\top\delta_{ij})
\]
by minimizing the cross-entropy against the \emph{soft responsibilities}
\(\gamma_{ij}(\ell)\) (equivalently, weighted logistic regression), with ridge
regularization by default and Firth penalization for small local samples.
\why{when clusters separate cleanly in \(\delta\)-space --- the good case --- the unpenalized
MLE diverges; and hard labels overstate the information available.}
Then \(\widehat\nu_i=\widehat\beta_i/\|\widehat\beta_i\|\) estimates the local boundary
normal and the intercept its signed displacement; sparse fits identify candidate gating
features; multiclass local fits serve junctions.

\subsection{Stage-repeating uncertainty}
\label{sec:gateuq}

Ordinary logistic standard errors ignore mixture uncertainty, synchronization error,
projection and scale selection, and label-assignment uncertainty. Gate (and composition, and
\(\widehat K\)) uncertainty is therefore reported from a \textbf{Stage-B bootstrap} that
resamples the recovery queries and repeats, per replicate: mixture fitting and
\(K_i^{\mathrm{loc}}\) selection, responsibility computation, partial matching,
synchronization, and gate fitting. Reported objects: normal cones for
\(\widehat\nu_i\), intervals for offsets and \(\widehat\pi_i\), persistence frequencies for
global classes.

\subsection{Certified propagation: unflagged is not interior}
\label{sec:certify}

An unflagged anchor may be truly interior --- or underpowered, probed in the wrong
directions, near a gate with overlapping outputs, part of a rare regime, or a multiplicity
miss. Non-rejection is not evidence that no boundary exists, so labels are \emph{not}
propagated through unflagged regions by default; doing so would confidently fill exactly the
regions where detection lacked sensitivity. Instead, each graph edge used for propagation
must be \textbf{certified} by bridge probes:

\begin{enumerate}[label=\textbf{CERT\arabic*.},leftmargin=3.2em]
  \item Place bridge points along the segment (or manifold path) between neighboring anchors
  \(x_i,x_j\), and probe each with a small battery at a bridge scale.
  \item The edge is \emph{certified same-branch} if every bridge point passes the
  minimum-signal and dispersion checks with no modality evidence, and the unperturbed bridge
  responses interpolate within the posterior predictive tube of a single-branch fit joining
  the endpoint components.
  \item The edge is \emph{certified crossing} if an interior bridge point is flagged;
  bisection along the edge then localizes the boundary, feeding Stage C geometry.
  \item Otherwise the edge is \emph{uncertifiable}; any label carried across it is reported
  as \textsc{extrapolated}, severed from certified claims.
\end{enumerate}
Propagation abstains in graph components disconnected from any reliably labeled anchor and
wherever multiple global labels remain plausible. Every reported regime map distinguishes
\textsc{certified} from \textsc{extrapolated} territory.
\why{this replaces v1's ``propagate into unflagged interior regions,'' which treated absence
of evidence as evidence of absence.}

\section{Complete annotated algorithm}
\label{sec:algorithm}

\paragraph{Stage A: detection.}
\begin{enumerate}[label=\textbf{A\arabic*.},leftmargin=2.7em]
  \item \textbf{Design.} Fix \(\psi\), probe families, ladder, anchors; split the budget
  \(\mathcal S_1\)--\(\mathcal S_4\); register negative controls and support diagnostics.
  \emph{Every later claim is relative to these choices; registering them first prevents
  post-hoc tuning.}
  \item \textbf{Noise and whitening.} Estimate \(\widehat\Sigma_{\mathrm{obs}}\) from
  repeats (shrinkage if \(V\) large); whiten. \emph{Makes the noise floor isotropic so
  eigenvalue excess is interpretable.}
  \item \textbf{Dispersion curves.} Compute \(r_i(\sigma_t)\) on \(\mathcal S_1\); apply the
  minimum-signal rule. \emph{Separates ``flat/no information'' from ``null-consistent.''}
  \item \textbf{Null generator.} Fit the hierarchical GP \(\nullgen\) robustly
  (trimmed/calibration-split), with cross-fitting folds for testing. \emph{The null is a
  fitted population of smooth behaviors, so it must be estimated without contamination from
  the anchors being tested.}
  \item \textbf{Scale statistic and screening.} Compute \(T_i^{\mathrm{scale}}\) (or
  \(R_i^{\log}\)); simulate cross-fitted screening \(p\)-values; BH at rate \(q\) as a
  screening device to form \(F\). \emph{Concentrates the confirmation budget; carries no
  final FDR claim.}
  \item \textbf{Shape corroboration.} On \(\mathcal S_2\), dip test on the specified
  projections at \(\sigma_i^\star\). \emph{A bump alone is not evidence of populations;
  modality is the corroborating signature.}
  \item \textbf{Confirmation.} Omnibus selection-aware bootstrap for the existence claim;
  per-anchor \(\widehat p_i^{\mathrm{sel}}\) with dependence-robust correction for
  anchor-level claims. \emph{All adaptive choices are inside the replicate; calibration is
  conditional on \(\nullgen\).}
\end{enumerate}

\paragraph{Stage B: recovery (confirmed anchors only).}
\begin{enumerate}[label=\textbf{B\arabic*.},leftmargin=2.7em]
  \item \textbf{Recovery sample.} Fresh queries near \(\sigma_i^\star\) from
  \(\mathcal S_3\). \emph{Severs recovery from selection.}
  \item \textbf{Local mixtures.} Fit, select \(K_i^{\mathrm{loc}}\), retain
  responsibilities; soft compositions \(\widehat\pi_i\). \emph{Soft estimates are consistent
  under overlap and carry uncertainty forward.}
  \item \textbf{Anchor orientation.} Posterior-assign \(\tilde\psi_{i0}\).
  \emph{Identifies the branch owning the anchor without archetype assumptions.}
  \item \textbf{Synchronization.} SYNC1--7: partial matching with dummy cost, confidence
  forest, cycle audit, abstention, \(\widehat K\). \emph{Turns arbitrary local labels into
  global identities with an explicit failure signal.}
  \item \textbf{Gates.} Responsibility-weighted ridge/Firth logistic; sparse variant for
  gating features. \emph{Regularization survives clean separation; weights carry mixture
  uncertainty.}
  \item \textbf{Certification.} Bridge-probe certification of every propagation edge;
  bisection localizes crossings. \emph{Non-rejection at an anchor never licenses
  propagation through it.}
  \item \textbf{Propagation.} Labels flow only across certified same-branch edges; all else
  \textsc{extrapolated} or abstained. \emph{Certified and extrapolated territory are never
  mixed in one claim.}
  \item \textbf{Stage-B bootstrap.} Repeat B2--B7 per replicate for uncertainty on
  \(\widehat\pi_i\), \(\widehat\nu_i\), \(\widehat K\), and label persistence.
  \emph{Uncertainty must include every estimated layer, not only the last one.}
\end{enumerate}

\paragraph{Stage C: interpretation.}
\begin{enumerate}[label=\textbf{C\arabic*.},leftmargin=2.7em]
  \item \textbf{Mechanism characterization.} Features, thresholds, or boundary maps
  predicting recovered labels and normals, with uncertainty. \emph{Converts geometry into a
  candidate mechanism.}
  \item \textbf{Consequential association.} Association of routing with protected
  attributes, geography, outcomes; spatially aware inference. \emph{The audit question is
  who is affected, not only whether structure exists.}
  \item \textbf{Probe-support comparison.} Contrast on-manifold, ambient, and
  feature-targeted findings, stratified by support realism. \emph{Different supports see
  different edges of the same system.}
  \item \textbf{Report.} Existence, structure, interpretation stated separately with
  abstentions and controls. \emph{Later-stage failure does not erase earlier-stage
  evidence.}
\end{enumerate}

\section{Stage B assumptions and diagnostics}
\label{sec:assumptions}

Fix an anchor \(x_i\); write \(Y_{ij}=\psi(f(x_i+\delta_{ij}))\), hidden labels
\(A_{ij}=g(x_i+\delta_{ij})\), compositions \(\pi_i(k)\), and regime-conditioned laws
\(P_{ik}\). Recovery is not assumption-free; each assumption is paired with a diagnostic and
an abstention condition.

\paragraph{B1 (stationarity).} One stable routing law governs the audit:
\(A_{ij}=g(x_i+\delta_{ij})\), with any randomness stationary. Rules out a stateful
adversary recognizing audit traffic. \emph{Diagnostic:} interleave disguised and audit
queries; test batchwise drift; withhold static recovery under material drift.

\paragraph{B2 (probe accessibility).} Every recovered regime is reached by some probe family:
\(\exists\,i:\pi_i(k)\ge\pi_{\min}>0\). \(\widehat K\) counts observable, probe-accessible
regimes, not internal code paths. \emph{Diagnostic:} compare scales and supports; label
support-specific modes as such.

\paragraph{B3 (local distinguishability).} Locally present regimes differ in law:
\(P_{ik}\neq P_{i\ell}\), with bounded optimal local classification error
\(\inf_h\Prob(h(Y_{ij})\neq A_{ij}\mid x_i)\le\varepsilon_{\mathrm{loc}}\). No Gaussian
assumption is intrinsic. \emph{Diagnostic:} held-out likelihood, stability, separation
statistics; decline decomposition when unstable.

\paragraph{B4 (prevalence).} \(m\,\pi_i(k)\ge n_{\min}\) (soft mass). Rare components are
not practically recoverable. \emph{Diagnostic:} report soft masses and CIs; merge or omit
failing components.

\paragraph{B5 (within-regime continuity).}
\(d_{\mathcal P}(P_{ik},P_{jk})\le L\,d_{\mathcal X}(x_i,x_j)\) for neighbors on one branch;
permits drifting \(w_k(x)\), forbids nothing about global variation. \emph{Diagnostic:}
smoothness of matched centroids/covariances along the graph.

\paragraph{B6 (matching margin).}
\(d_{\mathcal P}(P_{ik},P_{jk})+\gamma<\min_{\ell\neq k}d_{\mathcal P}(P_{ik},P_{j\ell})\),
\(\gamma>0\); otherwise branches can cross and exchange identities unobserved.
\emph{Diagnostic:} report best-vs-second-best margins, match-stability bootstrap,
\(\rho_{cc}\).

\paragraph{B7 (overlap connectivity).} Local views must chain (e.g.\ \(\{1,2\},\{2,3\},
\{3,4\}\)); the anchor--regime bipartite graph must be connected over each recovered family.
\emph{Diagnostic:} connectivity and articulation points; deletion stability; disconnected
families reported separately.

\paragraph{B8 (anchor fidelity).} \(A_{i0}=g(x_i)\): the unperturbed query is routed by the
same mechanism. Used only for orientation. \emph{Diagnostic:} repeated unperturbed queries
consistently join one component.

\paragraph{B9 (gate regularity).} Near a binary boundary, a local linear approximation
\(\beta_{i0}+\beta_i^\top\delta=0\) is adequate at the probed scale; sparsity is an
additional assumption used only when claiming a gating feature. \emph{Diagnostic:} held-out
label prediction; linear-vs-nonlinear comparison; spatial coherence of normals.

\paragraph{B10 (transversality).} \(\nu_i^\top\Sigma_{\mathrm{probe}}\nu_i>0\); probes
parallel to the boundary see nothing. \emph{Diagnostic:} compare supports; negative
conclusions stay scale- and support-relative.

\paragraph{Identifiability statement.} Under B1--B10 the target is recovery of observable
regimes up to global label permutation. Semantic names (``biased,'' ``innocuous'') are
produced by Stage C, never by clustering.

\section{Stage C: interpretation and audit finding}
\label{sec:stagec}

Multiple regimes do not imply misconduct; benign mixtures of experts, safety mechanisms, and
operational controllers are legitimately gated. Stage C characterizes the mechanism
(features, thresholds, boundary maps with uncertainty), tests consequential association
(protected attributes, geography, energy burden, service quality, error rates --- with
spatially aware inference), and compares probe supports with deliberately weakened language:
ambient-only evidence is \emph{consistent with} an off-manifold scaffold; on-manifold
evidence indicates behavior relevant to the modeled data support; a claim of misconduct
requires a separate interpretation and impact analysis beyond either. The reporting contract
separates \textbf{existence} (calibrated Stage-A evidence), \textbf{structure} (regimes,
compositions, partition, gates, with uncertainty, certification status, and abstentions), and
\textbf{interpretation} (the argued connection to a consequential or adversarial mechanism).
Failure at a later stage does not erase an earlier result.

\section{Validation and abstention contract}
\label{sec:abstention}

Recovery is withheld or qualified when any of the following fail: minimum-signal condition;
calibrated screening and omnibus confirmation; cluster stability under bootstrap and held-out
prediction; robustness of \(K_i^{\mathrm{loc}}\) to tuning; component separation relative to
within-component covariance; matching margins and cycle consistency \(\rho_{cc}\);
cross-support compatibility of shared components; held-out reconstruction of labels and
responses; edge certification along every propagation path; negative-control sanity (controls
must not be flagged at rates above their nominal levels); support realism thresholds for any
finding claimed to be operationally meaningful.

Reported quantities: omnibus and per-anchor selection-adjusted \(p\)-values;
insufficient-signal, abstention, and certification rates; cluster stability; component
separation; \(\rho_{cc}\); held-out assignment error; soft-composition intervals; normal
cones; \(\widehat K\) persistence; negative-control outcomes; support realism distributions.

\section{Relationship to separable profile-mean recovery}
\label{sec:spa}

The profile-mean approach factors \(M=WH+N\) and identifies simplex vertices from pure or
nearly pure columns (SPA with simplex-constrained regression). It requires separability,
constant mode signatures \(w_k\), and survives on profile means --- averaging away exactly
the within-profile multimodality that Stage B consumes. The present route replaces
separability with B1--B10, allows \(w_k(x)\) to drift (B5), and moves the difficulty from
vertex identification to local mixture estimation plus partial-permutation synchronization.
Neither dominates universally; when both structures hold, SPA provides a cross-check, and the
evaluation plan includes it as a baseline.

\section{Evaluation plan}
\label{sec:eval}

\paragraph{Suites.}
(1) smooth single-regime nulls with varying slope and curvature;
(2) heteroskedastic and anisotropic-noise single-regime nulls;
(3) two-regime threshold gates with controlled jump, noise, spatial dependence;
(4) three-regime Slack-style stacks with separate on-manifold and ambient probes;
(5) overlapping regimes approaching the distinguishability limit;
(6) disconnected coverage and deliberately inconsistent matchings;
(7) \textbf{misspecification}: heavy-tailed noise, unmodeled heteroskedasticity, kinks,
wrong Mat\'ern \(\nu\), adversarial smoothing near probe scales;
(8) \textbf{adaptive adversary probes}: response smoothing, noise inflation, probe-pattern
detection (scoped; full treatment deferred);
(9) baselines: SPA profile-mean recovery, global clustering, residual-outlier auditing,
detector-boundary anomaly flagging.

\paragraph{Metrics.}
Detection: realized FDR and power against jump size and noise; calibration curves of nominal
vs.\ realized error under correct and misspecified generators; stability comparison of
\(R_i\) (v1 ratio) vs.\ \(R_i^{\log}\) vs.\ \(T_i^{\mathrm{scale}}\); negative-control flag
rates. Recovery: query-label accuracy up to permutation; \(\|\widehat\pi_i-\pi_i\|\);
partition accuracy on certified territory; \(\widehat K\) distribution; boundary-normal
angular error; certification precision (rate at which certified same-branch edges truly
contain no crossing). Abstention: \textbf{operating characteristics} --- abstention rate in
non-identifiable settings vs.\ identifiable ones; the method should abstain where it should
and rarely where it should not.

\section{Bringing Stage B to Stage A's level}
\label{sec:parity}

Stage A now has a precise null claim, a concrete calibration route, a stabilized statistic,
and a falsification plan. Stage B reaches the same standard when it has:

\begin{enumerate}[label=\textbf{P\arabic*.},leftmargin=2.6em]
  \item \textbf{A formal estimand and loss.} The target: the restriction of the partition
  \(\{g^{-1}(k)\}\) to probe-accessible territory, up to label permutation; the composition
  field \(\pi_i\); local normals \(\nu_i\). Losses: permutation-minimized partition error on
  certified territory, \(\|\widehat\pi_i-\pi_i\|\), angular normal error. Stage A has a
  hypothesis; Stage B needs an estimand of the same precision.
  \item \textbf{Fully committed procedures.} This note commits the synchronization
  (SYNC1--7) and certification (CERT1--4) algorithms; remaining freedom --- \(\chi\),
  confidence thresholds, \(n_{\min}\), bridge densities --- needs registered defaults with
  sensitivity analyses, so that two auditors given the same data produce the same output.
  \item \textbf{Uncertainty that propagates every layer.} The Stage-B bootstrap of
  \S\ref{sec:gateuq} must cover mixture fitting, model selection, matching,
  synchronization, and gate fitting jointly; reported intervals that skip a layer are
  understatements.
  \item \textbf{A theorem-shaped target.} Minimal version: under B3--B7 with margin
  \(\gamma\), matching-cost concentration, and per-edge certified-matching error
  \(\alpha_e\), forest propagation recovers global labels up to permutation with probability
  at least \(1-\sum_{e\in\text{forest}}\alpha_e\); local clustering error enters
  \(\alpha_e\) through existing finite-sample mixture bounds; soft compositions satisfy a
  CLT under the local model. None of this is claimed proven here; it is the target that
  makes Stage B a result rather than a pipeline.
  \item \textbf{Abstention operating characteristics.} Quantitative evidence, from suite
  (6) and the misspecification suite, that the diagnostics abstain when assumptions fail
  --- ``knows when it does not know'' as a measured property, not a design intention.
  \item \textbf{Inherited selection discipline.} Stage B runs on selected anchors;
  estimates reusing selection data inherit winner's-curse bias. The budget split
  (\(\mathcal S_3\), \(\mathcal S_4\)) is the design-level fix; any reuse must enter the
  Stage-B bootstrap.
\end{enumerate}

\section{Claim boundary}
\label{sec:boundary}

This note specifies a complete, revised algorithmic proposal. The central research risk is no
longer whether the recovery idea makes sense; it is whether Stage A can make a precise,
calibrated claim that distinguishes boundary-like latent behavior from ordinary smooth
complexity, heteroskedasticity, and null-model misspecification. That is addressed by the
concrete null generator, cross-fitted and simulation-calibrated inference, the stabilized
scale statistic, and the falsification suites --- and it must be demonstrated empirically
before Stage A is presented as a result. Stage B is presented honestly as a partially
specified structural-recovery layer whose outputs are reported only where synchronization,
certification, and gate diagnostics pass, with \S\ref{sec:parity} as its maturation
contract. Detection remains the most developed component; recovery is the proposed next
stage; interpretation is a substantive argument, not an algorithmic output.

\end{document}
